\documentclass[UTF8,11pt,fontset=none]{ctexart}
\usepackage[a4paper,margin=1in]{geometry}
\usepackage{fontspec}
\setCJKmainfont{Songti SC}
\setCJKsansfont{PingFang SC}
\setCJKmonofont{Kaiti SC}
\usepackage{amsmath,amssymb,amsthm,mathtools}
\usepackage{booktabs,longtable,array}
\usepackage{enumitem}
\usepackage{xcolor}
\usepackage[colorlinks=true,linkcolor=blue,citecolor=blue,urlcolor=blue]{hyperref}

\setlist{nosep}
\setlength{\parskip}{0.65em}
\setlength{\parindent}{0pt}

\newtheorem{definition}{定义}[section]
\newtheorem{proposition}[definition]{命题}
\newtheorem{remark}[definition]{说明}
\newtheorem{example}[definition]{例子}
\newtheorem{principle}[definition]{原则}

\newcommand{\R}{\mathbb{R}}
\newcommand{\N}{\mathbb{N}}
\newcommand{\Vocab}{\mathcal V}
\newcommand{\D}{\mathcal D}
\newcommand{\X}{\mathcal X}
\newcommand{\Y}{\mathcal Y}
\newcommand{\ThetaSet}{\Theta}
\newcommand{\Loss}{\mathcal L}
\newcommand{\E}{\mathbb E}
\newcommand{\softmax}{\operatorname{softmax}}
\newcommand{\LN}{\operatorname{LN}}
\newcommand{\Attn}{\operatorname{Attn}}
\newcommand{\MHAttn}{\operatorname{MHA}}
\newcommand{\FFN}{\operatorname{FFN}}
\newcommand{\onehot}{\operatorname{onehot}}
\newcommand{\argmin}{\operatorname*{arg\,min}}
\newcommand{\argmax}{\operatorname*{arg\,max}}

\title{给范畴论背景读者的 AI 算法讲义：神经网络、Transformer 与 LLM}
\author{}
\date{\today}

\begin{document}
\maketitle

\begin{abstract}
本文面向具有范畴论或 topos 背景、但不熟悉现代深度学习算法细节的读者。本文的目的不是讨论工程调参，而是建立一套统一符号，用以描述监督学习、神经网络、反向传播、卷积网络、图神经网络、Transformer 和大语言模型中的基本数学对象。若要进一步讨论这些对象是否能被组织成范畴、预层、层、stack 或 topos，首先需要明确 AI 侧的对象、态射候选、局部结构和训练目标。
\end{abstract}

\section{目的与基本记号}

从范畴论观点出发，神经网络首先可以被视为一个映射
\[
f:X\to Y.
\]
这种观点刻画了固定参数后的输入--输出行为，但不足以描述机器学习中实际使用的模型结构。一个可训练模型通常同时包含以下四类数据：
\begin{enumerate}
\item 输入、输出和标签所在的数据空间；
\item 参数空间；
\item 用于评价预测误差的损失函数；
\item 在参数空间中更新参数的训练算法。
\end{enumerate}
因此，本文采用的基本记号是
\[
N_\theta:\X\to\Y,\qquad \theta\in\ThetaSet,
\]
其中 \(\ThetaSet\) 表示参数空间。训练过程并不是研究单个固定映射 \(N_\theta\)，而是在参数空间 \(\ThetaSet\) 中寻找使经验损失较小的参数。

\begin{principle}[三个相互区分的对象]
讨论 AI 与范畴论的关系时，至少应区分以下三类对象：
\begin{enumerate}
\item 模型族：\(\{N_\theta:\X\to\Y\}_{\theta\in\ThetaSet}\)；
\item 训练过程：\(\theta_0,\theta_1,\ldots,\theta_T\)；
\item 训练后模型：固定参数 \(\theta_T\) 后得到的函数 \(N_{\theta_T}\)。
\end{enumerate}
若一个范畴论表述没有说明它指向的是模型族、训练过程还是训练后模型，那么该表述的数学对象尚未被充分确定。
\end{principle}

\section{监督学习的统一符号}

\begin{definition}[数据集]
监督学习数据集记为
\[
\D=\{(x_i,y_i)\}_{i=1}^n,\qquad x_i\in\X,\quad y_i\in\Y.
\]
这里 \(\X\) 是输入空间，\(\Y\) 是标签或输出空间。
\end{definition}

典型例子如下：

\begin{center}
\begin{tabular}{lll}
\toprule
任务 & 输入 \(x\) & 输出 \(y\) \\
\midrule
图像分类 & 图像张量 \(x\in\R^{H\times W\times C}\) & 类别 \(y\in\{1,\ldots,K\}\) \\
回归 & 特征向量 \(x\in\R^d\) & 实数或向量 \(y\in\R^m\) \\
语言模型 & token 序列 \(x=(t_1,\ldots,t_m)\) & 下一个 token \(t_{m+1}\) \\
序列到序列 & 输入序列 \(x\) & 输出序列 \(y\) \\
\bottomrule
\end{tabular}
\end{center}

\begin{definition}[损失函数]
模型 \(N_\theta:\X\to\Y'\) 的损失函数通常写作
\[
\ell(N_\theta(x),y)\in\R_{\ge 0}.
\]
经验风险是训练集上的平均损失：
\[
\Loss_{\D}(\theta)=\frac1n\sum_{i=1}^n \ell(N_\theta(x_i),y_i).
\]
训练的基本目标是
\[
\theta^\ast\approx \argmin_{\theta\in\ThetaSet}\Loss_{\D}(\theta).
\]
\end{definition}

注意 \(\Y'\) 不一定等于 \(\Y\)。分类任务中，标签 \(y\) 是类别，而模型输出通常是 logits：
\[
z=N_\theta(x)\in\R^K.
\]
经过 softmax 得到概率分布
\[
p_k=\softmax(z)_k=\frac{\exp(z_k)}{\sum_{j=1}^K\exp(z_j)}.
\]
若真实类别为 \(y\)，交叉熵损失为
\[
\ell(z,y)=-\log p_y.
\]

\section{全连接神经网络}

最基本的神经网络是若干仿射映射和非线性函数的复合。令
\[
h_0=x\in\R^{d_0}.
\]
第 \(l\) 层定义为
\[
a_l=W_lh_{l-1}+b_l,\qquad h_l=\sigma_l(a_l),
\]
其中
\[
W_l\in\R^{d_l\times d_{l-1}},\qquad b_l\in\R^{d_l}.
\]
参数为
\[
\theta=(W_1,b_1,\ldots,W_L,b_L).
\]
输出为
\[
N_\theta(x)=h_L.
\]

\begin{remark}
从范畴论眼光看，层复合
\[
N_\theta=f_{L,\theta_L}\circ\cdots\circ f_{1,\theta_1}
\]
像是态射复合。但机器学习还关心每个 \(f_{l,\theta_l}\) 的参数如何变化、损失如何沿复合求导、以及哪些参数共享或受约束。这些信息不在普通态射复合记号中自动出现。
\end{remark}

常用非线性包括
\[
\operatorname{ReLU}(u)=\max(u,0),
\qquad
\tanh(u)=\frac{e^u-e^{-u}}{e^u+e^{-u}}.
\]
非线性函数的作用是改变可表示函数类。若没有非线性，多个线性层的复合仍是一个线性层：
\[
W_L\cdots W_1x.
\]
于是深度不会增加表达能力。

\section{训练与反向传播}

训练通常使用随机梯度下降及其变体。给定 mini-batch
\[
B\subseteq\{1,\ldots,n\},
\]
mini-batch 损失为
\[
\Loss_B(\theta)=\frac1{|B|}\sum_{i\in B}\ell(N_\theta(x_i),y_i).
\]
最简单的梯度下降更新是
\[
\theta_{t+1}=\theta_t-\eta_t\nabla_\theta\Loss_B(\theta_t),
\]
其中 \(\eta_t>0\) 是学习率。

\subsection{反向传播与链式法则}

反向传播是链式法则在层状复合函数中的系统化计算。设单个样本的损失为
\[
\Loss=\ell(h_L,y).
\]
定义误差信号
\[
\delta_l=\frac{\partial \Loss}{\partial a_l}\in\R^{d_l}.
\]
最后一层由损失函数给出 \(\delta_L\)。对中间层有递推
\[
\delta_l=(W_{l+1}^\top\delta_{l+1})\odot \sigma_l'(a_l),
\]
其中 \(\odot\) 是逐坐标乘法。于是参数梯度为
\[
\frac{\partial \Loss}{\partial W_l}=\delta_l h_{l-1}^\top,
\qquad
\frac{\partial \Loss}{\partial b_l}=\delta_l.
\]

\begin{principle}[反向传播的范畴论读法]
若前向传播是复合
\[
f_L\circ\cdots\circ f_1,
\]
则反向传播沿相反方向传播导数信息。它不是把原函数反过来，而是把 cotangent/gradient 信息通过每一层的导数转置向后传递。
\end{principle}

\section{卷积网络：局部性和权重共享}

卷积神经网络 CNN 主要用于图像和局部网格数据。输入可写成
\[
x\in\R^{H\times W\times C_{\mathrm{in}}}.
\]
卷积核为
\[
K\in\R^{r\times s\times C_{\mathrm{in}}\times C_{\mathrm{out}}}.
\]
输出特征图 \(y\in\R^{H'\times W'\times C_{\mathrm{out}}}\) 的坐标为
\[
y_{u,v,c}
=
\sum_{\alpha=1}^r\sum_{\beta=1}^s\sum_{c'=1}^{C_{\mathrm{in}}}
K_{\alpha,\beta,c',c}\,
x_{u+\alpha,v+\beta,c'}.
\]

这里有两个关键结构。

第一是局部性：\(y_{u,v,c}\) 只依赖输入中靠近 \((u,v)\) 的一个小窗口。第二是权重共享：同一个卷积核 \(K\) 在所有位置使用。这使 CNN 比全连接网络更适合表达平移相关的图像结构。

如果要用层或 sheaf 语言看 CNN，图像 patch 可以作为局部上下文；patch 之间的重叠给出相容性问题；卷积核共享则是一种对平移结构的约束。

\section{图神经网络：邻域聚合}

图神经网络 GNN 处理图
\[
G=(V,E).
\]
每个节点 \(v\in V\) 有表示
\[
h_v^{(0)}\in\R^{d_0}.
\]
第 \(l\) 层通常由消息传递给出：
\[
m_v^{(l)}
=
\operatorname{AGG}^{(l)}
\{\,\phi^{(l)}(h_v^{(l)},h_u^{(l)},e_{uv})\mid u\in \mathcal N(v)\,\},
\]
\[
h_v^{(l+1)}
=
\psi^{(l)}(h_v^{(l)},m_v^{(l)}).
\]
这里 \(\mathcal N(v)\) 是 \(v\) 的邻居集合，\(\operatorname{AGG}\) 通常是求和、平均或最大值。

GNN 的基本结构要求之一是置换等变性。若重新编号节点，输出表示也应按同样方式重新编号，而不应依赖任意编号。这一性质与范畴论中的自然性和等变性有相似之处；但在具体模型中，它需要由聚合函数对邻居顺序的不敏感性来保证。

\section{Transformer 的统一符号}

Transformer 是现代 LLM 的基础架构之一。设词表为
\[
\Vocab=\{1,\ldots,V\}.
\]
一个长度为 \(T\) 的 token 序列写作
\[
t_{1:T}=(t_1,\ldots,t_T),\qquad t_i\in\Vocab.
\]
嵌入矩阵为
\[
E\in\R^{V\times d}.
\]
token \(t_i\) 的向量表示是
\[
x_i=E_{t_i}+p_i\in\R^d,
\]
其中 \(p_i\) 是位置编码。把所有位置堆叠成矩阵
\[
X=
\begin{bmatrix}
x_1^\top\\
\vdots\\
x_T^\top
\end{bmatrix}
\in\R^{T\times d}.
\]

\subsection{Scaled dot-product attention}

给定
\[
X\in\R^{T\times d},
\]
定义 query、key、value：
\[
Q=XW_Q,\qquad K=XW_K,\qquad V=XW_V,
\]
其中
\[
W_Q,W_K\in\R^{d\times d_k},
\qquad
W_V\in\R^{d\times d_v}.
\]
attention 输出为
\[
\Attn(Q,K,V)
=
\softmax\!\left(\frac{QK^\top}{\sqrt{d_k}}+M\right)V.
\]
矩阵
\[
\frac{QK^\top}{\sqrt{d_k}}\in\R^{T\times T}
\]
给出每个位置对每个位置的相似度。softmax 按行计算，因此第 \(i\) 行是位置 \(i\) 对所有位置 \(j\) 的权重。

在自回归语言模型中，mask \(M\) 用来禁止看未来 token：
\[
M_{ij}=
\begin{cases}
0,& j\le i,\\
-\infty,& j>i.
\end{cases}
\]
这保证预测第 \(i+1\) 个 token 时只能使用 \(t_1,\ldots,t_i\)。

\subsection{Multi-head attention}

多头注意力把 \(H\) 个 attention head 并行计算：
\[
\operatorname{head}_r
=
\Attn(XW_Q^{(r)},XW_K^{(r)},XW_V^{(r)}),
\qquad r=1,\ldots,H.
\]
然后拼接并线性变换：
\[
\MHAttn(X)=
\operatorname{Concat}(\operatorname{head}_1,\ldots,\operatorname{head}_H)W_O.
\]
多头机制允许不同 head 学习不同关系，例如局部语法、长程依赖、实体指代或格式模式。这里的“关系”不是人工指定的逻辑关系，而是由训练数据和损失函数诱导出来的统计结构。

\subsection{Transformer block}

一个常见 decoder-only Transformer block 可写为
\[
Y=X+\MHAttn(\LN(X)),
\]
\[
Z=Y+\FFN(\LN(Y)).
\]
其中 FFN 是逐位置的两层 MLP：
\[
\FFN(u)=W_2\sigma(W_1u+b_1)+b_2.
\]
注意 attention 在 token 位置之间混合信息，而 FFN 对每个位置分别作用。

\section{语言模型目标}

语言模型学习条件概率
\[
p_\theta(t_{1:T})=\prod_{i=1}^T p_\theta(t_i\mid t_{<i}).
\]
训练目标通常是最大化训练文本的似然，等价于最小化负对数似然：
\[
\Loss(\theta)
=
-\sum_{i=1}^T \log p_\theta(t_i\mid t_{<i}).
\]
模型在每个位置输出 logits
\[
z_i\in\R^{|\Vocab|}.
\]
概率为
\[
p_\theta(t_{i+1}=v\mid t_{\le i})
=
\softmax(z_i)_v.
\]

\begin{remark}
从数学上看，LLM 不是直接输出“意义”，而是输出下一个 token 的条件分布。意义、推理、工具使用和知识检索是从大规模训练、架构约束、上下文输入和解码策略中涌现或被工程化出来的行为。
\end{remark}

\section{推理、采样与 KV cache}

训练后，给定 prompt
\[
c=(t_1,\ldots,t_m),
\]
模型逐步生成：
\[
t_{m+1}\sim p_\theta(\,\cdot\mid t_{\le m}),
\]
\[
t_{m+2}\sim p_\theta(\,\cdot\mid t_{\le m+1}),
\]
依此类推。

常见解码方式包括：
\[
t_{i+1}=\argmax_{v\in\Vocab}p_\theta(v\mid t_{\le i})
\]
的 greedy decoding，以及按概率采样的 stochastic decoding。temperature \(\tau>0\) 调整分布尖锐程度：
\[
p_\tau(v)=
\softmax(z/\tau)_v.
\]
\(\tau\) 越小，输出越确定；\(\tau\) 越大，输出越随机。

KV cache 是推理加速结构。由于每一步都会重复使用过去 token 的 key 和 value，模型把已经计算的
\[
K_{\le i},\quad V_{\le i}
\]
缓存起来。生成第 \(i+1\) 个 token 时，只需为新 token 计算新的 \(q,k,v\)，再与缓存交互。这不改变模型定义，只改变计算效率。

\section{RAG 与工具调用}

检索增强生成 RAG 把外部文档引入上下文。给定查询 \(q\)，检索器从文档库
\[
\mathcal R=\{d_1,\ldots,d_M\}
\]
取出相关文档
\[
r_1,\ldots,r_k.
\]
然后模型条件化于扩展上下文：
\[
p_\theta(t_{i+1}\mid q,r_1,\ldots,r_k,t_{\le i}).
\]
这不是改变 Transformer 的基本公式，而是改变输入上下文和信息来源。

工具调用可抽象为模型在某一步输出一个结构化 action
\[
a_i\in\mathcal A,
\]
外部环境返回观察
\[
o_i\in\mathcal O.
\]
之后模型继续在扩展历史
\[
(t_{\le i},a_i,o_i)
\]
上生成。若要用范畴论语言描述工具调用，态射不应只表示文本变换，还要表示模型、工具、环境之间的信息流。

\section{这些 AI 对象怎样和范畴/topos 问题对接}

现在可以更精确地说 AI 侧有哪些候选结构。

\begin{longtable}{p{0.28\textwidth}p{0.64\textwidth}}
\toprule
AI 对象 & 可能的范畴论/topos 对接点 \\
\midrule
模型族 \(\{N_\theta\}_{\theta\in\ThetaSet}\) & 参数化态射、函数对象、模型空间 \\
训练轨迹 \(\theta_0,\ldots,\theta_T\) & 优化动力系统，不等同于固定态射 \\
层 \(f_{l,\theta_l}\) & 复合结构、局部模块、可替换组件 \\
CNN patch & 局部上下文、覆盖、限制映射 \\
GNN 邻域 & 局部邻域聚合、置换等变性 \\
attention 矩阵 & token 间上下文依赖，不是简单局部限制 \\
LLM prompt & 条件上下文，可作为某类 context object \\
RAG 文档 & 外部证据片段，可能形成覆盖或索引数据 \\
工具调用 & 模型与环境的交互态射或过程 \\
\bottomrule
\end{longtable}

\begin{principle}
若要把一个 AI 系统放入 topos 框架，不能只说“存在上下文”。必须说明上下文对象是什么、态射是什么、覆盖是什么、局部数据是什么、相容性是什么、粘合后得到什么。
\end{principle}

\section{思维链的谨慎解释}

Chain-of-thought 可以看成中间文本轨迹
\[
c=(s_1,\ldots,s_m),
\]
其中每个 \(s_i\) 是一个中间推理片段。但这些片段不一定彼此一致。模型可能先提出错误假设，再修正；也可能并列比较多个方向。

因此，用 sheaf 的“局部相容则全局粘合”直接解释 chain-of-thought 会过于简单。更合理的做法是把思维链看成一种动态过程：
\[
\text{state}_0\to \text{state}_1\to\cdots\to \text{state}_m,
\]
其中态射可能表示扩展、反驳、修正、压缩或验证。若要建立范畴模型，应先决定这些态射的类型，而不是预设它们都是限制映射。

\section{最小词汇表}

\begin{center}
\begin{tabular}{ll}
\toprule
AI 术语 & 本讲义符号 \\
\midrule
输入空间 & \(\X\) \\
输出/标签空间 & \(\Y\) \\
参数空间 & \(\ThetaSet\) \\
模型 & \(N_\theta:\X\to\Y'\) \\
数据集 & \(\D=\{(x_i,y_i)\}_{i=1}^n\) \\
损失 & \(\Loss_{\D}(\theta)\) \\
隐藏状态 & \(h_l\) 或 \(X\in\R^{T\times d}\) \\
logits & \(z\in\R^K\) 或 \(z_i\in\R^{|\Vocab|}\) \\
token 序列 & \(t_{1:T}\) \\
语言模型概率 & \(p_\theta(t_i\mid t_{<i})\) \\
attention & \(\Attn(Q,K,V)\) \\
\bottomrule
\end{tabular}
\end{center}

\section{结语}

范畴论读者理解 AI 算法时，最重要的是不要把所有东西都压成一个函数。固定参数后的网络确实是函数；但机器学习真正关心的是参数化函数族、训练目标、优化路径、局部模块、上下文依赖、概率输出和交互过程。只有这些对象被清楚区分之后，讨论预层、层、stack 或 topos 才有可靠的落点。

\end{document}
